\item \textbf{Problema de repaso: Enfriamiento y atrapamiento de átomos por láser.}
Steven Chu, Claude Cohen-Tannoudji y William Phillips recibieron el Premio Nobel de Física en 1997 por sus métodos para enfriar átomos con luz láser. Un haz de átomos de masa $m \approx 10^{-25}\text{ kg}$ se mueve a $1\text{ km/s}$. Un láser sintonizado a una transición de $500\text{ nm}$ se dirige en sentido opuesto. El átomo en estado fundamental absorbe un fotón. Después de $\approx 10^{-8}\text{ s}$ el átomo sufre emisión espontánea aleatoria en cualquier dirección, por lo que el retroceso promedio de la emisión es cero.
\begin{enumerate}
    \item Calcule la desaceleración promedio que experimenta el haz atómico debido a los sucesivos ciclos de absorción y retroceso.
    \item Estime la distancia geométrica que requiere el haz para detenerse por completo.
\end{enumerate}